% arara: pdflatexmk

\documentclass[%
	shownavigationbar=	false,%
	showsidebar= 		false,%
	fleqn,%
	xcolor=			{usenames, svgnames},%
	usepdftitle=			false,%
%	compress,%
%	draft,%
]{beamer}

\usepackage{%
	marvosym,%
	graphicx,%
	hyperref,%
	parskip,%
	booktabs,%
	tabularx,%
	multicol,%
	microtype,%
%	etex,%
	array,%
	colortbl,%
	xspace,%
	enumerate,%
	amsmath,%
}

\definecolor{VUBlue1}{cmyk}{1,0,0,0}
\definecolor{VUBlue2}{cmyk}{1,0.16,0,0.20}
\definecolor{VUBlue3}{cmyk}{1,0.55,0,0.18}
\definecolor{FEWEB1}{cmyk}{0.60,0.90,0,0}
\definecolor{FEWEB2}{cmyk}{0.5,1,0,0}
\definecolor{FEWEB3}{cmyk}{0.30,0.90,0,0}
\definecolor{Ivory2}{rgb}{0.93,0.93,0.88}
\definecolor{Ivory3}{rgb}{0.8,0.8,0.76}

\usepackage{tikz}
\usetikzlibrary{positioning, backgrounds, fit, arrows.meta, calc}% fit, backgrounds,arrows.meta,

\newcommand{\link}[2]{\tikz[baseline=-2pt]\node[rectangle, inner sep=2pt, rounded corners=.75pt, draw=VUBlue3] {\href{#1}{\texttt{#2}}};\xspace}

\newcommand{\bula}{{\tikz\node[rectangle, inner sep=2pt, rounded corners=.75pt, draw=VUBlue3, thick, fill=VUBlue3] {};}\xspace}%fill=VUBlue1!75
\newcommand{\bulb}{\tikz\node[rectangle, inner sep=2pt, rounded corners=.75pt, draw=VUBlue3, thick] {};\xspace}
\newcommand{\bulc}{\tikz\draw[VUBlue3, thick] (0,0) (0,.3ex)--(5pt,.3ex) (.6em,0);\xspace}%(0,.75ex)--(.6em,.75ex)

\usepackage{enumitem}%
\setlist[itemize,1]{label={\raisebox{0.2ex}\bula}, itemsep=.5ex, topsep=1ex}%
% LEVEL 2
\setlist[itemize,2]{label={\raisebox{0.2ex}\bulb}, itemsep=.5ex, topsep=.75ex}%
% LEVEL 3
\setlist[itemize,3]{label=\bulc, itemsep=.25ex, topsep=.5ex}%

\usepackage[T1]{fontenc}

\usepackage[UKenglish]{babel}

\usetheme{default}
\mode<presentation>
\usecolortheme{default}
\usecolortheme{rose}
\setbeamercovered{transparent=30}

\setbeamercolor{structure}{fg=VUBlue3}%titlelike
\setbeamercolor{normal text}{fg=black}
\setbeamercolor{alerted text}{fg=FEWEB2}

\setbeamertemplate{subsection in toc}{%
	\leavevmode\leftskip=5.65ex%
	\llap{\raisebox{0.3ex}{{\bulc}}\kern1ex}%
  	\inserttocsubsection\par%
}
\setbeamertemplate{subsection in toc shaded}[default][0]%
\setbeamertemplate{title page}[default][left]%
\setbeamertemplate{navigation symbols}{}%
\setbeamertemplate{blocks}[rounded][shadow=false]
\setbeamercolor{block title}{bg=VUBlue2, fg=white}%
\setbeamercolor{block body}{bg=VUBlue2!5, fg=black}%
\setbeamercolor{block title alerted}{bg=FEWEB1, fg=white}%
\setbeamercolor{block body alerted}{bg=FEWEB2!10, fg=black}%
\newcommand{\plus}{\textcolor{Green}{$+$}\ }
\newcommand{\minus}{\textcolor{Red}{$-$}\ }

%%%%
\title{Finding information efficiently}
\author{\large Matthijs de Zwaan\\}%\url{m.c.de.zwaan@vu.nl}

\date{28 October 2016}

\begin{document}

% TITLE PAGE
\begin{frame}[plain]
{\vskip 3ex \Large\inserttitle}
\begin{columns}[t]
\column{.2\textwidth}
\column{.8\textwidth}
\vskip -5ex
{\large\insertauthor}
\vskip 1ex
\insertinstitute
\vskip 5ex
\insertdate
\end{columns}
%\vspace(3ex)
%\centering
\begin{tikzpicture}[remember picture,overlay]
%	\node {\includegraphics[width=\paperwidth]{UB_ENG_Wit}};
	\node[anchor=south, at={(current page.base)}] {\includegraphics[width=\paperwidth]{UB_ENG_Blauw_FC}};
\end{tikzpicture}
\end{frame}


\begin{frame}
\begin{tikzpicture}[remember picture,overlay]
	\node[at={(current page.center)}] {\includegraphics[trim=.5cm .5cm .5cm .5cm, clip,  width=\paperwidth]{hobbyhorse}};
\end{tikzpicture}
\end{frame}

\section*{Copy, combine, create}

\begin{frame}
\frametitle{Copy, combine, create (and don't be shy about it!)}

\begin{block}{\link{http://quoteinvestigator.com/2013/03/06/artists-steal/}{``Good artists copy; great artists steal.''}}
Academics steal and leave finger prints.
\end{block}
\vspace{1ex}
\begin{block}{Accountability is important:}
\begin{itemize}
	\item What did you copy?
	\item What is your own work?
\end{itemize}
\end{block}
\vspace{1ex}
\begin{alertblock}{Search for information takes time}
Two hours is not enough!
\end{alertblock}
\end{frame}

% Object/Preview
\begin{frame}<1,2>[label=preview]
\frametitle{\inserttitle}
\begin{columns}[t]
\column{.2\textwidth}
\column{.8\textwidth}
\vskip 2ex
\uncover<1,2>{
	{\only{\color{VUBlue3}}<2>Finding literature efficiently}\\
	\textcolor{black!60}{A workflow for finding literature}
}
\vskip 3ex
\uncover<1,3>{
	{\only{\color{VUBlue3}}<3>Finding economic data efficiently}\\
	\textcolor{black!60}{Data sources}\\
	\textcolor{black!60}{Research data services}
}
\end{columns}
\end{frame}

\begin{frame}<1>[label=workflow]
\frametitle{A workflow for finding literature}
\begin{tikzpicture}[
	auto,%
	-{Latex[round]},%
	node distance=.5cm and .75cm,%
	shift={(current page.center)},%
	shorten <=2pt, shorten >=2pt,%
	remember picture, overlay,%
	very thick,%
	looseness=1.4,%
]
\node[rectangle, rounded corners, draw] (topic) at (-4.5cm, 3cm) {Topic search};
\node[rectangle, rounded corners, draw, below right={of topic.south}, text width=7em] (sel1) {Select, scan \&\\select again};
\node[rectangle, rounded corners, draw, below right={of sel1.south}] (snow) {Snowball search};
\node[rectangle, rounded corners, draw, below right={of snow.south}] (sel2) {Select};
\node[rectangle, rounded corners, draw, below right={of sel2.south}] (read) {Read \& select};
\node[rectangle, rounded corners, draw, below right={of read.south}] (cite) {Reference};

\draw (topic.south) to [out=-120,in=180] (sel1.west);
\draw (sel1.south) to [out=-120,in=180] (snow.west);
\draw (snow.south) to [out=-120, in=180] (sel2.west);
\draw (sel2.south) to [out=-120, in=180] (read.west);
\draw (read.south) to [out=-120, in=180] (cite.west);

\draw[dashed] (sel1.north) to [out=60,in=0] (topic.east);
\draw[dashed] (read.north) to [out=60,in=0] (snow.east);
\draw[dashed, -, shorten >=3pt] (read.north) to [out=60,in=0, curve to] node[text width=4em] {Repeat if necessary} (topic.east);

\only<2->{
\begin{scope}[on background layer, rounded corners]
	\node[fill=VUBlue2!15,fit=(topic)(sel1)] {};
\end{scope}
}
\only<3->{
\begin{scope}[on background layer, rounded corners]
	\node[fill=VUBlue2!15,fit=(snow)(sel2)] {};
\end{scope}
}
\end{tikzpicture}
\end{frame}

\section{Finding (and using) literature efficiently}
\subsection{Google versus the rest}

\begin{frame}
\frametitle{Google versus the rest}
\begin{block}{GIYF, sometimes}
\uncover<1>{
\begin{itemize}
	\item Use \link{https://scholar.google.nl/}{Google Scholar}, not ``plain'' Google.
	\item Easiest option if you know what you're looking for.
	\item Not for books! Use \link{https://vu.on.worldcat.org/advancedsearch?databaseList=}{LibSearch}.
\end{itemize}
}
\end{block}

\begin{block}{The rest: bibliographic databases}
\uncover<2>{
\begin{itemize}
	\item More structured search
	\item Tools for expanding/limiting your search
	\item Searches in more specific (economic/business) database
\end{itemize}
}
\end{block}
\end{frame}

\subsection{Topic search}

\begin{frame}
\frametitle{Topic search}
\begin{block}{Search using keywords}
\begin{itemize}
	\item operators: `AND' or `OR' (see \link{http://images.webofknowledge.com/WOK46/help/WOS/ht_operators.html}{here} or \link{https://youtu.be/5K_yUMvt_nU}{here}),\\
	\begin{columns}
	\column{.1\linewidth}
	\column{.4\linewidth}
	\bulb $x\text{ AND }y$:
	\begin{tikzpicture}[baseline=-1mm]
		\begin{scope}[thick]
			\draw[clip] (0,0) circle[radius=3mm];
			\node at (2.5mm,0) [fill=VUBlue2, shape=circle, minimum size=6mm, inner sep=0pt] {};
		\end{scope}
		\node at (2.5mm,0) [shape=circle, minimum size=6mm, inner sep=0pt, draw=black] {};
		\node at (-1.4mm,0){$ \scriptstyle x$};
		\node at (3.9mm,0){$\scriptstyle y$};
	\end{tikzpicture}% (0.3,0) circle (.5cm);
	\column{.4\linewidth}
	\bulb $x\text{ OR }y$: \tikz[baseline=-1mm] \draw[fill=VUBlue2] (0,0) circle[radius=3mm] (2.5mm,0) circle[radius=3mm];% (0.3,0) circle (.5cm);
	\column{.1\linewidth}
	\end{columns}
	\item Also useful:\\
	\begin{tikzpicture}[anchor=west]
		\node at (1em,0) {\bulb parentheses: (\ldots)};
		\node at (1em,-3ex) {\bulb quotes: \textquotedbl\ldots\textquotedbl};
		\node at (13em,0) {\bulb wildcards: `*', `?'};
	\end{tikzpicture}
	\item (fdi OR "foreign direct investment" OR multinational) AND (tfp OR "total factor productivity")
\end{itemize}
\end{block}
\begin{block}{Bibliographic databases very powerful}
\begin{itemize}
	\item \link{http://www.ubvu.vu.nl/pub/index_oclc.cfm?SearchObjectId=7&objectid=109&ordering=1&openitem=1288}{Web of Science}: general database, widely used
	\item \link{http://www.ubvu.vu.nl/pub/index_oclc.cfm?SearchObjectId=7&objectid=109&ordering=1&openitem=182}{ABI/Inform}: business/economics, thesaurus
	\item \link{http://www.ubvu.vu.nl/pub/index_oclc.cfm?SearchObjectId=7&objectid=109&ordering=1&openitem=2055}{Business Source Elite}: business/economics, thesaurus
	\item Also use other filters: time, field of study, etc
\end{itemize}
\end{block}
\end{frame}

\begin{frame}
\frametitle{Good key words are key}
\begin{block}{Where to start}
\begin{itemize}
	\item Research question gives you keywords:\\
	``How does $X$ affect $Y$?''\\
	\item Textbook, article, Wikipedia\\
	Textbook: check the index \\
	Articles: often keywords listed below abstract
\end{itemize}
\end{block}
\begin{block}{Where to find alternatives}
	\begin{itemize}
	\item Include different spelling (US/GB: s-z, o-ou); \\
	Use wild cards: `internationali?ation', `internat*'
	\item Thesaurus (next slides)
	\item J of Economic Literature codes (next slides)
\end{itemize}
\end{block}
%\end{columns}
\end{frame}

\begin{frame}[plain]
\frametitle{\link{http://zbw.eu/stw/descriptor/10823-6}{STW thesaurus for economics}}
\begin{tikzpicture}
	\node[at=(current page.center)] {\includegraphics[height=\textheight]{zbw}};
\end{tikzpicture}
\end{frame}

\begin{frame}[plain]
\frametitle{\link{https://www.aeaweb.org/jel/guide/jel.php}{JEL codes}}
\begin{tikzpicture}
	\node at (current page.center) {\includegraphics[width=.95\paperwidth]{JEL}};
\end{tikzpicture}
\end{frame}

\begin{frame}[plain]
\frametitle{Web of Science: example}
\begin{tikzpicture}
	\node[at=(current page.center)] {\includegraphics[width=.95\paperwidth]{WoSsearch}};
\end{tikzpicture}
\end{frame}

\subsection{Selecting articles}

\begin{frame}
\frametitle{How to select articles?}
\begin{block}{Rules of thumb}
\begin{itemize}
%	\item Newspaper articles, blogs etc: good for inspiration or background info, not as `serious' reference \\(but cite them if used!)
	\item Peer-reviewed work (journal articles) preferred over non-peer-reviewed (eg, working papers).
	\item (Exceptions: \link{http://www.nber.org/papers.html}{NBER}, \link{http://www.iza.org/en/webcontent/publications/papers}{IZA}, \ldots)
	\item Articles with more citations are often more influential
%	\item Ask your professor (\textit{but first do your homework}).
	\item No guarantee! Use your brain first, then your advisor.
\end{itemize}
\end{block}

\begin{block}{For your paper (if you need $X$ articles)}
\begin{itemize}
	\item Do topic search $\Rightarrow$ \textit{many} articles
	\item Sort by citations
	\item Select approx $1.5\times X$ top articles
	\item Scan / read superficially (abstract/intro/results)
	\item Discard what isn't useful
	\item Repeat \& adjust search if necessary
\end{itemize}
\end{block}
\end{frame}

\subsection{Snowball search}

\againframe<2>{workflow}

\begin{frame}
\frametitle{How to search: Forward and backward snowball}
\begin{block}{Backward snowball}
\begin{itemize}
	\item Which articles were used for this research?
	\item Check the references in an article
\end{itemize}
\end{block}

\begin{block}{Forward snowball}
\begin{itemize}
	\item How was this article used in other research?
	\item Google Scholar: ``Cited by: \ldots''
    \item WoS: ``Times cited: \ldots''
\end{itemize}
\end{block}
\end{frame}

\begin{frame}[plain]
    \begin{tikzpicture}[remember picture,overlay]
		\node[at=(current page.center)] (gs1a) {
			\includegraphics[width=\paperwidth]{googlescholar1a}
		};
		\only<2>{
		\node[at=(current page.center)] {
			\includegraphics[width=.95\paperwidth]{googlescholar1b}
		};
		\path[draw, color=red, thick] (0,0.3)--(0,1.3)--(6,1.3)--(6,0.3)--cycle;
		}
	\end{tikzpicture}
\end{frame}

\begin{frame}[plain]
    \begin{tikzpicture}[remember picture,overlay]
		\node[at=(current page.center)] {
			\includegraphics[width=\paperwidth, trim=4cm 0 4.5cm 0, clip]{WoS1a}
		};
		\only<2>{
		\node[at=(current page.center), anchor=north] {
			\includegraphics[width=.8\paperwidth]{WoS1b}
		};
		\path[draw, color=red, thick] (2,-0.75)--(2,0.3)--(10.5,0.3)--(10.5,-0.75)--cycle;
		}
	\end{tikzpicture}
\end{frame}

\begin{frame}
\begin{block}{For your paper}
\begin{itemize}
	\item Select 5--10 of most interesting papers
	\item Backward snowball to see if you missed anything
	\item Forward snowball: \\
	Do any recent papers cite many of your selected papers?
	\item If yes: scan for interest. \\
	New theories, new data, new controversy?
\end{itemize}
\end{block}
\end{frame}

\againframe<3>{workflow}

\subsection{Proper references}

\begin{frame}
\frametitle{Leave your finger prints}
\begin{block}{Proper references}
\begin{itemize}
	\item Plagiarism is \alert{serious} academic crime!
	\item Most common in Econ \& Finance: author-year format
	\item author(s) and year (in parenthesis) of publication in text
	\item Full reference in bibliography\\(At least: authors, year of pub, full article title, journal title, journal volume/issue, pages.  Include a \link{https://en.wikipedia.org/wiki/Digital_object_identifier}{doi} if you can.)
	\item See \link{http://libguides.vu.nl/B_business_admin/incorporating_literature}{here} for more
\end{itemize}
\end{block}

\begin{block}{Invest some time in (semi-)automation}
\begin{itemize}
	\item \link{http://endnote.com/}{Endnote}: supported by \link{http://www.ub.vu.nl/en/help-advice/tools/endnote/index.aspx}{library}.
	\item \link{https://www.mendeley.com/}{Mendeley}, \link{https://www.zotero.org/}{Zotero}: online \& free
	\item \link{http://www.jabref.org/}{JabRef}, \link{http://bibdesk.sourceforge.net/}{BibDesk}: \LaTeX{} users
	\item Ready-made formats, integration with word processors.
\end{itemize}
\end{block}
\end{frame}

\begin{frame}[plain]
\begin{tikzpicture}[remember picture,overlay]
	\node[at=(current page.center)] {
		\includegraphics[width=\paperwidth]{WoS2a}
	};
	\only<2>{
	\node[at=(current page.center), anchor=south west] (zoom1) {
		\includegraphics[width=.4\paperwidth]{WoS2b}
	};
	\path[draw, color=red, thick] (3.1,1) --+ (0,1)--+(2.25,1)--+(2.25,0)--cycle;%, name path={box1}
	\draw[color=red, thick, ->] (4.2,1) to [bend right] ($(zoom1.west)+(0,-.5)$);

	\node[at=(current page.west), anchor=north west] (zoom2) {
		\includegraphics[width=.6\paperwidth]{WoS2c}
	};
	\path[draw, color=red, thick] (-.5,-.75) --+ (0,1.25)--+(2.75,1.25)--+(2.75,0)--cycle;%, name path=box2
	\draw[color=red, thick, ->] (2.25,-.25) to [out=0, in=90] (zoom2.north);
	}
\end{tikzpicture}
\end{frame}

\section{Finding (and using) economic data efficiently}
\againframe<3>{preview}

\begin{frame}
\frametitle{Useful secondary data}
\begin{block}{Freely available on the web}
\begin{multicols}{2}
\begin{itemize}
	\item \link{http://ec.europa.eu/eurostat/data/database}{Eurostat}
	\item \link{https://fred.stlouisfed.org/}{Federal reserve (US)}
	\item \link{http://data.worldbank.org/}{World bank}
	\item \link{https://data.oecd.org/}{OECD}
	\item \link{https://www.imf.org/en/Data}{IMF}
	\item Recently: data available with publications.
	\item and \link{http://rfe.org/showCat.php?cat_id=2}{many}, \link{http://www.ubvu.vu.nl/pub/index_oclc.cfm?SearchObjectId=7&aantalitem=30&lang=&fromitem=1&max=77&ordering=1&objectid=109&trefwoord=&fctN=st3&fctIncV=13&fctExcV=+&fctN=st5&fctIncV=17&fctExcV=+}{many} others\ldots
\end{itemize}
\end{multicols}
\end{block}
\end{frame}

\begin{frame}
\frametitle{Research Data Services}
\begin{block}{The library manages and supports many data subscriptions}
\begin{itemize}
	\item Listed on the UB Blackboard page
	\item Firm level financial data: \\Amadeus (EU), Zephyr (M\&A), Compustat (US \& Global)
	\item Market data: CRSP, NYSE TaQ
	\item Datastream
\end{itemize}
\end{block}
\begin{block}{Help desk services}
\begin{itemize}
	\item Mail for short questions, appointments for more difficult problems
	\item Workshops for finding/using data, and for merging several different data sets
	\item Support for Excel and Stata (data management, not analysis!)
\end{itemize}
\end{block}
\end{frame}

\begin{frame}[plain]
\frametitle{Research Data Services}
\begin{tikzpicture}
	\node[at=(current page.west), anchor=west] {\includegraphics[width=\paperwidth]{BBa}};
	\only<2>{
		\path[draw, color=red, thick] (0,3) --+ (0,4)--+(2,4)--+(2,0)--cycle;%, name path=box1
		\node[at=(current page.center),anchor=east] {\includegraphics[height=.75\paperheight]{BBb}}	;
	}
\end{tikzpicture}
\end{frame}

\begin{frame}[plain]
\begin{columns}[t]
\column{.2\textwidth}
\column{.8\textwidth}
Questions?\\[2ex]
Send us an email, make an appointment.
\end{columns}

\begin{tikzpicture}[remember picture,overlay]
%	\node {\includegraphics[width=\paperwidth]{UB_ENG_Wit}};
	\node[anchor=south, at={(current page.base)}] {\includegraphics[width=\paperwidth]{UB_ENG_Blauw_FC}};
\end{tikzpicture}
\end{frame}

\begin{frame}%[allowframebreaks]
\frametitle{About the library (1)}
\begin{block}{The physical Library}
\begin{itemize}
	\item Many \link{http://ub.vu.nl/en/facilities/index.aspx}{facilities}, first floor of main building.
\end{itemize}
\end{block}

\begin{block}{The digital Library}
\begin{itemize}
	\item \link{http://ub.vu.nl/en/index.aspx}{http://ub.vu.nl/} for information, opening hours, news
\end{itemize}
\end{block}
\end{frame}

\begin{frame}
\frametitle{About the library (2)}
\begin{block}{Research Data Services}
\begin{itemize}
	\item \link{http://ub.vu.nl/en/education-research/research-data-services/support-for-secondary-sources/index.aspx}{Home page} for opening hours, mail, reservations for data computers, etc
	\item Help desk (walk-in or email), guidelines, online tutorials
	\item Register with the RDS Blackboard page for an overview of available data sources, training sessions, etc
	\begin{itemize}
		\item Log in at \link{https://bb.vu.nl/}{Blackboard}
		\item Courses \textgreater ``Universiteitsbibliotheek VU'' \textgreater ``UBVU\_DATACENTRUMUBVU\_9999\_1'' \textgreater enroll
	\end{itemize}
\end{itemize}
\end{block}
\end{frame}

\begin{frame}
\begin{tikzpicture}[remember picture,overlay]
	\node[at={(current page.center)}] {\includegraphics[width=\paperwidth]{cowboy}};
\end{tikzpicture}
\end{frame}

\end{document}
